%% =====================================================================
%%  Unit 16 — My plans  (Parler de ses projets)
%%  Pistes : 16.01–16.04 dialogues · 16.05–16.06 vocabulaire · 16.07–16.08 grammaire
%%           16.09 prononciation · 16.10–16.13 exercices (textes en annexe)
%% =====================================================================

\newcommand{\otIllusUnXVI}{\otIllusDuo{trip/104/86/15/0.10,summer/190/92/20/0.09,vacation/122/14/13/0.08,
  soon/150/8/11/0.08,travel/196/22/12/0.07,next/100/50/12/0.07}{emma}{What are you going to do?}{lea}{I'm going to travel!}}

% Carte schématique de la Route 66 (exercice 6)
\newcommand{\otRouteSoixanteSix}{%
  \begin{tikzpicture}[x=1mm,y=1mm,font=\fontsize{8.4}{10}\selectfont]
    \fill[otNavyTint,rounded corners=2mm] (-4,-4) rectangle (162,40);
    \fill[otNavy!12] (-4,-4) -- (-4,40) -- (2,40) .. controls (4,28) and (0,14) .. (3,-4) -- cycle;
    \node[rotate=90,text=otNavy!60,font=\otSemi\fontsize{6.4}{8}\selectfont] at (-0.8,18) {\enb{PACIFIC OCEAN}};
    \draw[otOrange,line width=2.2pt,line cap=round,line join=round]
      (8,10) .. controls (18,12) and (24,16) .. (32,18) .. controls (42,20) and (48,15) .. (55,15)
      .. controls (64,15) and (70,13) .. (78,13) .. controls (86,13) and (92,11) .. (98,12)
      .. controls (110,14) and (120,17) .. (128,20) .. controls (136,23) and (144,28) .. (150,30);
    \foreach \x/\y/\v/\a in {8/10/Los Angeles/below,32/18/Flagstaff/below,55/15/Albuquerque/below,
      78/13/Amarillo/below,98/12/Oklahoma City/below,128/20/St.\ Louis/below,150/30/Chicago/above}{%
      \fill[white,draw=otOrangeDark,line width=0.9pt] (\x,\y) circle (1.5);
      \node[\a=1.6mm,text=otInk] at (\x,\y) {\enb{\v}};}
    \otDessinIcone{etoile}{otGreenDark}{otNavyTint}{6mm}{(36,31)}
    \node[right,text=otGreenDark,font=\otSemi\fontsize{8}{10}\selectfont] at (39,31) {\enb{Grand Canyon}};
    \draw[otGreenDark,dashed,line width=0.5pt] (35,28.5) -- (32.6,19.8);
    \node[draw=otInk,fill=white,line width=0.8pt,rounded corners=1mm,inner sep=1.2mm,
      font=\otBold\fontsize{9}{10}\selectfont] at (113,33) {\enb{ROUTE 66}};
    \node[text=otMuted,font=\fontsize{7.4}{9}\selectfont] at (113,26.5) {\enb{2,448 miles · 3,940 km}};
  \end{tikzpicture}}

\unite{16}{My plans}{Parler de ses projets}{otIllusUnXVI}

\situation{Épisode 16}{Mi-décembre. À l'approche des fêtes, chacun parle de ses projets au studio :
Emma rentre à Londres, et les parents de Léa vont venir la voir à New York ! Le soir, au café, Jack
raconte à Sam le grand voyage qu'il prépare pour l'été.}

\begin{objectifs}
\item parler de vos projets et de vos intentions avec \en{be going to} ;
\item poser des questions sur les projets des autres ;
\item situer une action dans le futur : \en{tomorrow}, \en{next week}, \en{in two weeks}… ;
\item parler de voyages et de vacances ;
\item comprendre l'anglais parlé : \en{gonna}, \en{wanna}.
\end{objectifs}

\auprogramme
  {Parler de ses projets et de ses intentions}
  {\en{be going to}\par\vspace{0.6mm}Les expressions du futur}
  {Voyages, vacances, projets : 24 expressions}
  {\en{gonna}, \en{wanna} : l'anglais parlé}
  {Le \en{road trip} américain}

\semaine{casque/otNavy/Dialogues/Écouter et lire les deux scènes,
  bulle/otGreenDark/Vocabulaire/Apprendre les 24 expressions,
  blocs/otViolet/Grammaire/be going to et le futur,
  sons/otOrangeDark/Sons et culture/gonna et le road trip,
  crayon/otNavy/Exercices/12 exercices corrigés,
  coche/otGreenDark/Bilan/L'essentiel et le test express}

\clearpage
%% ---------------------------------------------------------------------
\rubrique{ecouter}{Projets de vacances}

\begin{dialogue}{1}{What are you going to do?}{Vendredi 12 décembre, au studio. Dernière réunion avant
  les vacances de Noël.}{16.01}{16.02}
\dit{emma}{So, Léa, \hl{what are you going to do} for Christmas?\note{1}}{Alors, Léa, qu'est-ce que tu vas faire pour Noël ?}
\dit{lea}{\hl{My parents are going to visit} me! They arrive on December twentieth.}{Mes parents vont venir me voir ! Ils arrivent le 20 décembre.}
\dit{emma}{That's wonderful! \hl{Are you going to show} them the city?}{C'est merveilleux ! Tu vas leur faire visiter la ville ?}
\dit{lea}{Yes! We're going to see the big Christmas tree, and \hl{we're going to go} ice skating.}{Oui ! On va voir le grand sapin de Noël, et on va faire du patin à glace.}
\dit{jack}{And you, Emma? \hl{Are you going to fly} to London?}{Et toi, Emma ? Tu vas prendre l'avion pour Londres ?}
\dit{emma}{\hl{Yes, I am}. I'm going to spend two weeks with my family.\note{2}}{Oui. Je vais passer deux semaines avec ma famille.}
\dit{lea}{What about you, Jack?}{Et toi, Jack ?}
\dit{jack}{\hl{I'm not going to travel}. I'm going to stay home… and sleep!\note{3}}{Je ne vais pas voyager. Je vais rester chez moi… et dormir !}
\end{dialogue}

\begin{comprendre}
\pointnote{1} \en{\textbf{What are you going to do?}} \en{be going to} + base verbale exprime un
  \textbf{projet}, une intention : c'est le futur proche du français (« je vais… »).
\pointnote{2} \en{\textbf{Yes, I am.}} La réponse courte reprend \en{be} : \en{Yes, I am} / \en{No,
  I'm not}.
\pointnote{3} \en{\textbf{I'm not going to travel.}} Négation : \en{be} + \en{not} + \en{going to} +
  base verbale.
\end{comprendre}

\begin{dialogue}{2}{A road trip}{Le soir, au café où travaille Sam. Jack a étalé une carte des
  États-Unis sur le comptoir.}{16.03}{16.04}
\dit{sam}{Hey, that's a big map! What's the plan?}{Hé, c'est une grande carte, ça ! C'est quoi, le projet ?}
\dit{jack}{\hl{Next summer}, I'm going to drive across America, on Route 66!\note{4}}{L'été prochain, je vais traverser l'Amérique en voiture, par la Route 66 !}
\dit{sam}{From Chicago to Los Angeles? That's crazy! \hl{How long is it going to take}?}{De Chicago à Los Angeles ? C'est fou ! Ça va prendre combien de temps ?}
\dit{jack}{About three weeks. \hl{I'm gonna rent} a car and sleep in motels.\note{5}}{Environ trois semaines. Je vais louer une voiture et dormir dans des motels.}
\dit{sam}{Are you going to see the Grand Canyon?}{Tu vas voir le Grand Canyon ?}
\dit{jack}{Of course! It's \hl{on the way}.}{Bien sûr ! C'est sur la route.}
\dit{sam}{Hey, can I come with you? \hl{I wanna see} the desert!}{Hé, je peux venir avec toi ? Je veux voir le désert !}
\dit{jack}{Sure! But you're going to pay for the gas!}{Bien sûr ! Mais c'est toi qui vas payer l'essence !}
\end{dialogue}

\begin{comprendre}
\pointnote{4} \en{\textbf{Next summer}} Les expressions du futur se placent en début ou en fin de
  phrase : \en{next summer}, \en{next week}, \en{tomorrow}. Attention : pas de \en{the} devant \en{next}.
\pointnote{5} \en{\textbf{I'm gonna rent…}} À l'oral, \en{going to} se prononce souvent \en{gonna}, et
  \en{want to}, \en{wanna}. On les entend partout, mais on ne les écrit pas (sauf dans un texto).
\end{comprendre}

%% ---------------------------------------------------------------------
\rubrique{mots}{Les mots de l'unité}

Voici les 24 expressions de l'unité. En américain, les vacances se disent \en{vacation} ; au
Royaume-Uni, \en{holiday}.

\begin{vocabulaire}{Voyages et vacances}{16.05}
\categorie{Voyager}
\mot{a trip}{ə ˈtrɪp}{un voyage}
\mot{travel}{ˈtrævəl}{voyager}
\mot{a vacation}{ə veɪˈkeɪʃən}{des vacances \registre{américain}}
\mot{visit}{ˈvɪzɪt}{rendre visite à ; visiter}
\mot{stay}{steɪ}{rester, séjourner}
\mot{fly}{flaɪ}{prendre l'avion, voler}
\mot{a flight}{ə ˈflaɪt}{un vol}
\mot{rent a car}{ˌrent ə ˈkɑːr}{louer une voiture}
\categorie{Préparer son voyage}
\mot{a suitcase}{ə ˈsuːtkeɪs}{une valise}
\mot{pack}{pæk}{faire ses bagages}
\mot{the beach}{ðə ˈbiːtʃ}{la plage}
\mot{the mountains}{ðə ˈmaʊntənz}{la montagne}
\mot{go camping}{ˌɡoʊ ˈkæmpɪŋ}{faire du camping}
\end{vocabulaire}

\begin{vocabulaire}{Parler de l'avenir}{16.06}
\categorie{Quand ?}
\mot{tomorrow}{təˈmɑːroʊ}{demain}
\mot{this weekend}{ðɪs ˈwiːkend}{ce week-end}
\mot{next month}{ˌnekst ˈmʌnθ}{le mois prochain}
\mot{next year}{ˌnekst ˈjɪr}{l'année prochaine}
\mot{in two weeks}{ɪn ˌtuː ˈwiːks}{dans deux semaines}
\mot{soon}{suːn}{bientôt}
\mot{one day}{ˌwʌn ˈdeɪ}{un jour}
\categorie{Les projets}
\mot{a plan}{ə ˈplæn}{un projet}
\mot{What are you going to do?}{ˌwʌt ɑːr juː ˌɡoʊɪŋ tə ˈduː}{Qu'est-ce que tu vas faire ?}
\mot{I'm going to travel.}{aɪm ˌɡoʊɪŋ tə ˈtrævəl}{Je vais voyager.}
\mot{I'm not going to work.}{aɪm ˌnɑːt ˌɡoʊɪŋ tə ˈwɝːk}{Je ne vais pas travailler.}
\end{vocabulaire}

%% ---------------------------------------------------------------------
\rubrique[be going to]{grammaire}{\en{be going to} : parler de ses projets}

\otIntertitre[otViolet]{1 · \en{be going to} + base verbale}
\begin{center}\otTableau
\begin{tabular}{@{}l@{\hspace{8mm}}>{\enOn}l<{\enOff}@{\hspace{8mm}}l@{}}
\otTH{Forme} & \otTH{Exemple} & \otTH{Français}\\\midrule
affirmation      & I'\textbf{m going to} travel. / She'\textbf{s going to} travel. & je vais voyager…\\
négation         & I'\textbf{m not going to} work. / They \textbf{aren't going to} work. & je ne vais pas…\\
question         & \textbf{Are} you \textbf{going to} fly? / What \textbf{is} he \textbf{going to} do? & tu vas… ?\\
réponses courtes & Yes, I \textbf{am}. / No, she \textbf{isn't}. & oui / non\\\bottomrule
\end{tabular}\end{center}

\begin{regle}[Quand l'employer ?]
\begin{etapes}
\item Pour un \textbf{projet} déjà décidé : \en{I'm going to visit my parents next week.}
\item Pour une prévision évidente, quand on voit les signes : \en{Look at the clouds! It's going to rain.}
\item \en{be} change avec le sujet ; \en{going to} ne change jamais ; le verbe reste à la base verbale.
\end{etapes}
\end{regle}

\blocaudio{16.07}{Écoutez et répétez}{Six phrases avec \en{be going to}.}

\otIntertitre[otViolet]{2 · Les expressions du futur}
\begin{center}\otTableau
\begin{tabular}{@{}>{\enOn}l<{\enOff}@{\hspace{6mm}}l@{\hspace{14mm}}>{\enOn}l<{\enOff}@{\hspace{6mm}}l@{}}
\otTH{Anglais} & \otTH{Français} & \otTH{Anglais} & \otTH{Français}\\\midrule
tonight      & ce soir           & next week           & la semaine prochaine\\
tomorrow     & demain            & next summer         & l'été prochain\\
this weekend & ce week-end       & in two weeks        & dans deux semaines\\
soon         & bientôt           & one day             & un jour\\\bottomrule
\end{tabular}\end{center}

\blocaudio{16.08}{Écoutez et répétez}{Une question sur les projets, puis sa réponse.}

\begin{attention}[Les pièges des francophones]
\fauxjuste{I going to travel.}{I'm going to travel.}
\fauxjuste{I'm going travel.}{I'm going to travel.}
\fauxjuste{She's going to visits Boston.}{She's going to visit Boston.}
\fauxjuste{the next week}{next week}
\fauxjuste{I go to travel next summer.}{I'm going to travel next summer.}
\end{attention}

%% ---------------------------------------------------------------------
\rubrique[gonna et wanna]{sons}{\en{gonna}, \en{wanna} : l'anglais parlé}

Dans la conversation, les Américains « avalent » certains mots. \en{going to} devient \en{gonna}
\phon{ˈɡʌnə}, \en{want to} devient \en{wanna} \phon{ˈwɑːnə}, \en{got to} devient \en{gotta}
\phon{ˈɡɑːtə}. Ce n'est pas de l'anglais incorrect : c'est l'anglais parlé de tous les jours. Il
faut savoir le reconnaître ; à l'écrit, gardez toujours la forme complète.

\begin{phrases}{16.09}{Écoutez la forme parlée et répétez-la. La forme écrite est indiquée en gras.}
\phrase[I'm gonna travel.]{I'm going to travel.}{à l'oral : \en{I'm gonna travel}}
\phrase[What are you gonna do?]{What are you going to do?}{à l'oral : \en{What are you gonna do?}}
\phrase[We're gonna be late!]{We're going to be late!}{à l'oral : \en{We're gonna be late!}}
\phrase[I wanna go home.]{I want to go home.}{à l'oral : \en{I wanna go home}}
\phrase[Do you wanna come?]{Do you want to come?}{à l'oral : \en{Do you wanna come?}}
\phrase[I gotta go!]{I've got to go!}{à l'oral : \en{I gotta go!} (« je dois y aller »)}
\end{phrases}

\begin{attention}[À l'oral seulement]
\en{gonna} ne remplace \en{going to} que devant un verbe : \en{I'm gonna travel}. Quand on dit où l'on
va, on garde \en{going to} : \en{I'm going to London} (jamais \en{I'm gonna London}).
\end{attention}

%% ---------------------------------------------------------------------
\rubrique{culture}{Le road trip américain}

\begin{culture}[La Route 66]
Ouverte en 1926, la Route 66 reliait Chicago à Los Angeles sur 2\,448 miles, soit près de
4\,000\,km. Les Américains l'appellent \en{the Mother Road} ou \en{the Main Street of America}.
Remplacée par les autoroutes en 1985, elle reste le rêve des amateurs de \en{road trip} : motels aux
enseignes au néon, \en{diners} des années cinquante et paysages de désert.
\end{culture}

\begin{paire}
\begin{culture}[Les parcs nationaux]
Le premier parc national du monde, Yellowstone, a été créé aux États-Unis en 1872. Le pays en compte
aujourd'hui 63, dont le Grand Canyon, en Arizona, profond de plus de 1\,800 mètres.
\end{culture}
\begin{culture}[Miles et gallons]
Sur la route, les distances sont en \en{miles} (1 mile ≈ 1,6\,km) et la vitesse en \en{miles per
hour} : 65\,mph ≈ 105\,km/h. L'essence (\en{gas}) se vend au \en{gallon} (≈ 3,8 litres).
\end{culture}
\end{paire}

%% ---------------------------------------------------------------------
\rubrique{pratique}{Exercices}

\exercice[ecrire,niveau=1]{En voyage}{Écrivez l'expression anglaise.}
\begin{questions}[2][french]
\item un voyage → \enb{\trou[34mm]{a trip}}
\item louer une voiture → \enb{\trou[34mm]{rent a car}}
\item une valise → \enb{\trou[34mm]{a suitcase}}
\item prendre l'avion → \enb{\trou[34mm]{fly}}
\item des vacances (US) → \enb{\trou[34mm]{a vacation}}
\item faire ses bagages → \enb{\trou[34mm]{pack}}
\item la plage → \enb{\trou[34mm]{the beach}}
\item bientôt → \enb{\trou[34mm]{soon}}
\end{questions}

\exercice[ecrire,niveau=1]{\en{be going to}}{Complétez avec \en{be going to} et le verbe entre
parenthèses.}
\begin{anglais}
\begin{questions}
\item Emma \trou[46mm]{is going to fly} to London. \hfill (fly)
\item Léa's parents \trou[46mm]{are going to visit} New York. \hfill (visit)
\item We \trou[46mm]{are going to rent} a car. \hfill (rent)
\item Jack \trou[46mm]{isn't going to work} on Christmas Day. \hfill (not work)
\item They \trou[46mm]{aren't going to take} the bus. \hfill (not take)
\item \trou[46mm]{Are you going to travel} this summer? \hfill (you / travel)
\item What \trou[46mm]{is Sam going to do} tonight? \hfill (Sam / do)
\end{questions}
\end{anglais}

\exercice[lire,ecrire,niveau=2]{L'agenda de Léa}{Voici le programme de Léa pendant les vacances.
Répondez aux questions par une réponse courte ou une phrase complète.}
\begin{center}
\begin{tcolorbox}[enhanced,width=128mm,colback=white,colframe=otRule,boxrule=0.6pt,arc=2mm,
  left=4mm,right=4mm,top=3mm,bottom=3mm,
  title={\otSemi\fontsize{9}{11}\selectfont\enb{Léa's holiday plans}},coltitle=white,colbacktitle=otGreenDark,
  fonttitle=\otSemi,attach boxed title to top left={xshift=4mm,yshift=-2.5mm},boxed title style={arc=1mm,boxrule=0pt}]
\vspace{2mm}\renewcommand{\arraystretch}{1.45}
\begin{tabular}{@{}>{\otSemi\fontsize{9}{11}\selectfont\color{otGreenDark}\enOn}l<{\enOff}@{\hspace{6mm}}>{\enOn}l<{\enOff}@{}}
Sat., Dec. 20 & Mom and Dad arrive at JFK airport, 5 p.m.\\
Sun., Dec. 21 & Rockefeller Center: the Christmas tree + ice skating\\
Wed., Dec. 24 & Christmas Eve: dinner at Mrs.\ Lopez's apartment\\
Sat., Dec. 27 & The Museum of Modern Art with Jack's sister Kate\\
Wed., Dec. 31 & New Year's Eve in Times Square!\\
\end{tabular}
\end{tcolorbox}
\end{center}
\begin{anglais}
\begin{questions}
\item Are Léa's parents going to arrive on December 20th? \hfill \trou[50mm]{Yes, they are.}
\item Are they going to arrive in the morning? \hfill \trou[50mm]{No, they aren't.}
\item Is Léa going to go ice skating on December 24th? \hfill \trou[50mm]{No, she isn't.}
\item Where are they going to have dinner on Christmas Eve? \hfill \trou[50mm]{At Mrs. Lopez's apartment.}
\item Where are they going to be on December 31st? \hfill \trou[50mm]{In Times Square.}
\end{questions}
\end{anglais}

\exercice[ecrire,niveau=2]{Posez la question}{Écrivez la question avec \en{be going to}.}
\begin{questions}
\item \en{you} / \en{travel} / \en{this summer} \hfill \enb{\trou[76mm]{Are you going to travel this summer?}}
\item \en{Emma} / \en{fly} / \en{to London} \hfill \enb{\trou[76mm]{Is Emma going to fly to London?}}
\item \en{what} / \en{you} / \en{do} / \en{tomorrow} \hfill \enb{\trou[76mm]{What are you going to do tomorrow?}}
\item \en{where} / \en{they} / \en{stay} \hfill \enb{\trou[76mm]{Where are they going to stay?}}
\item \en{Jack} / \en{rent} / \en{a car} \hfill \enb{\trou[76mm]{Is Jack going to rent a car?}}
\end{questions}

\exercice[ecrire,niveau=2]{Dans l'ordre}{Remettez les mots dans l'ordre pour former une phrase.}
\begin{anglais}
\begin{questions}[1][english]
\item \etiq{going} \etiq{I'm} \etiq{travel} \etiq{to} \etiq{.}\hfill\trou[66mm]{I'm going to travel.}
\item \etiq{summer} \etiq{she's} \etiq{to} \etiq{next} \etiq{going} \etiq{fly} \etiq{.}\hfill\trou[66mm]{She's going to fly next summer.}
\item \etiq{not} \etiq{we're} \etiq{to} \etiq{going} \etiq{work} \etiq{.}\hfill\trou[66mm]{We're not going to work.}
\item \etiq{you} \etiq{going} \etiq{are} \etiq{stay} \etiq{to} \etiq{?}\hfill\trou[66mm]{Are you going to stay?}
\item \etiq{weeks} \etiq{in} \etiq{two} \etiq{see} \etiq{you} \etiq{!}\hfill\trou[66mm]{See you in two weeks!}
\end{questions}
\end{anglais}

\exercice[lire,niveau=2]{Le grand projet de Jack}{Jack envoie son itinéraire à Sam. Lisez le message et
la carte, puis indiquez si chaque phrase est vraie (V) ou fausse (F).}
\begin{center}\otRouteSoixanteSix\end{center}
\begin{anglais}
\begin{tcolorbox}[otBase,colback=white,colframe=otRule,boxrule=0.5pt,fontupper=\fontsize{10.6}{14.2}\selectfont]
Hi Sam! Here's my plan. I'm going to start in Chicago on June 15th. I'm going to drive about 300
miles a day, and I'm going to stop in St.\ Louis, Oklahoma City and Albuquerque. In Arizona, I'm going
to see the Grand Canyon. I'm not going to sleep in hotels: I'm going to stay in motels, and I'm going
to camp in the national parks. The trip is going to take three weeks, and I'm going to finish in Los
Angeles… on the beach! Are you going to come with me? — Jack
\end{tcolorbox}
\begin{questions}[2]
\item Jack is going to start in Los Angeles. \hfill\trou[9mm]{F}
\item He's going to drive about 300 miles a day. \hfill\trou[9mm]{V}
\item He's going to see the Grand Canyon. \hfill\trou[9mm]{V}
\item He's going to sleep in hotels. \hfill\trou[9mm]{F}
\item The trip is going to take three months. \hfill\trou[9mm]{F}
\item He's going to finish on the beach. \hfill\trou[9mm]{V}
\end{questions}
\end{anglais}

\exercice[ecouter,niveau=1,piste=16.10]{Maintenant ou plus tard ?}{Écoutez chaque phrase, dite deux
fois. Parle-t-on d'une action en cours (maintenant) ou d'un projet (plus tard) ? Cochez.}
\begin{questions}[3]
\item \choix{maintenant}\quad\choix{plus tard}
\item \choix{maintenant}\quad\choix{plus tard}
\item \choix{maintenant}\quad\choix{plus tard}
\item \choix{maintenant}\quad\choix{plus tard}
\item \choix{maintenant}\quad\choix{plus tard}
\item \choix{maintenant}\quad\choix{plus tard}
\end{questions}

\exercice[ecouter,ecrire,niveau=2,piste=16.11]{Les projets de chacun}{Quatre personnes parlent de
leurs projets pour les vacances. Associez chaque personne à son projet : écrivez la lettre.}
\begin{anglais}
\noindent\begin{minipage}[t]{0.40\linewidth}
\begin{questions}[1][english]
\item Emma \hfill\trou[9mm]{c}
\item Sam \hfill\trou[9mm]{a}
\item Mrs.\ Lopez \hfill\trou[9mm]{d}
\item Maya \hfill\trou[9mm]{b}
\end{questions}
\end{minipage}\hfill
\begin{minipage}[t]{0.54\linewidth}
\begin{questions}[1][english]
\item[a.] play concerts with a band
\item[b.] go skiing in Vermont
\item[c.] see her family in London
\item[d.] cook for twelve people
\item[e.] go to the beach in Florida
\end{questions}
\end{minipage}
\end{anglais}

\exercice[ecouter,ecrire,niveau=2,piste=16.12]{Dictée à trous}{Mrs.\ Lopez parle avec Léa dans
l'escalier. Écoutez et complétez.}
\begin{anglais}
\noindent\begin{tabular}{@{}l@{\hspace{6mm}}l@{}}
\otNomDit{lopez} & Léa, what are you \trou{going} to do for Christmas?\\[1.8mm]
\otNomDit{lea}   & My parents are going to \trou{visit} me!\\[1.8mm]
\otNomDit{lopez} & How nice! How \trou{long} are they going to stay?\\[1.8mm]
\otNomDit{lea}   & Two \trou{weeks}. They're going to \trou{fly} back on January third.\\[1.8mm]
\otNomDit{lopez} & Are they going to \trou{stay} with you?\\[1.8mm]
\otNomDit{lea}   & No, they \trou{aren't}. They're going to stay in a \trou{hotel}.\\
\end{tabular}
\end{anglais}

\exercice[ecrire,niveau=2]{Traduisez}{Traduisez en anglais.}
\begin{questions}[1][french]
\item Qu'est-ce que tu vas faire ce week-end ?\lignes{1}
\item Je vais rendre visite à mes parents.\lignes{1}
\item Nous n'allons pas prendre l'avion.\lignes{1}
\item Tu vas louer une voiture ? — Oui. \textit{(réponse courte)}\lignes{1}
\item L'été prochain, ils vont faire du camping.\lignes{1}
\item Elle va arriver dans deux semaines.\lignes{1}
\item Regarde ces nuages ! Il va pleuvoir.\lignes{1}
\item Un jour, je vais voir le Grand Canyon.\lignes{1}
\end{questions}

\exercice[ecrire,niveau=3]{Mes projets}{En six phrases, présentez vos projets pour l'année prochaine :
un voyage, un projet personnel, ce que vous n'allez pas faire.}
\lignes{6}

\exercice[ecouter,parler,niveau=3,piste=16.13]{À vous de parler !}{Emma vous interroge sur vos projets.
Répondez \textbf{à voix haute} pendant chaque pause.}
\begin{center}\otTableau
\begin{tabular}{@{}>{\otSemi\fontsize{8.6}{10}\selectfont\color{otNavy}}l@{\hspace{3mm}}>{\enOn}p{70mm}<{\enOff}@{\hspace{6mm}}p{78mm}@{}}
 & \otTH{Emma dit\ldots} & \otTH{Vous\ldots}\\\midrule
1. & What are you going to do this weekend? & dites deux choses que vous allez faire.\\
2. & Are you going to travel next summer? & répondez et dites où.\\
3. & How are you going to travel? & dites comment (avion, voiture, train).\\
4. & Are you going to work tomorrow? & répondez par une réponse courte.\\
5. & Great! Do you wanna have a coffee? & acceptez.\\\bottomrule
\end{tabular}\end{center}

%% ---------------------------------------------------------------------
\rubrique{bilan}{Faites le point}

\begin{essentiel}
\begin{minipage}[t]{0.54\linewidth}
\otEssTitre{Les phrases clés}
\renewcommand{\arraystretch}{1.32}
\begin{tabular}{@{}>{\enOn\otSemi\fontsize{9.4}{11}\selectfont}l<{\enOff}@{\hspace{3mm}}>{\fontsize{9.2}{11}\selectfont\color{white!78!otNavy}}l@{}}
What are you going to do?    & Qu'est-ce que tu vas faire ?\\
I'm going to travel.         & Je vais voyager.\\
I'm not going to work.       & Je ne vais pas travailler.\\
Are you going to fly?        & Tu vas prendre l'avion ?\\
Yes, I am. / No, I'm not.    & Oui. / Non.\\
See you next week!           & À la semaine prochaine !\\
\end{tabular}
\end{minipage}\hfill
\begin{minipage}[t]{0.41\linewidth}
\otEssTitre{\en{be going to}}
{\fontsize{9.2}{11.6}\selectfont \en{am / is / are} + \en{going to} + base verbale : projet, intention,
prévision évidente\par}
\otEssTitre{Le futur}
{\fontsize{9.2}{11.6}\selectfont \en{tomorrow} · \en{next week} (sans \en{the}) · \en{in two weeks} ·
\en{soon}\par}
\otEssTitre{À l'oral}
{\fontsize{9.2}{11.6}\selectfont \en{going to} → \en{gonna} · \en{want to} → \en{wanna}\par}
\end{minipage}
\end{essentiel}

\begin{jesais}
\item parler de mes projets avec \en{be going to} ;
\item poser des questions sur les projets des autres et y répondre ;
\item employer les expressions du futur (\en{tomorrow}, \en{next week}…) ;
\item parler de voyages et de vacances ;
\item reconnaître \en{gonna} et \en{wanna} à l'oral ;
\item présenter le \en{road trip} américain et la Route 66.
\end{jesais}

\clearpage
\begin{testexpress}
\qcm[c]{« Je vais voyager » se dit :}{\en{I go to travel.}}{\en{I going to travel.}}{\en{I'm going to travel.}}
\qcm[a]{« Tu vas louer une voiture ? » se dit :}{\en{Are you going to rent a car?}}{\en{Do you going to rent a car?}}{\en{You are going rent a car?}}
\qcm[b]{\en{Is she going to fly?} — Réponse négative :}{\en{No, she doesn't.}}{\en{No, she isn't.}}{\en{No, she not.}}
\qcm[b]{« la semaine prochaine » se dit :}{\en{the next week}}{\en{next week}}{\en{the week next}}
\qcm[c]{« dans deux semaines » se dit :}{\en{since two weeks}}{\en{two weeks ago}}{\en{in two weeks}}
\qcm[a]{« Il va pleuvoir » (on voit les nuages) se dit :}{\en{It's going to rain.}}{\en{It goes to rain.}}{\en{It's raining.}}
\qcm[c]{À l'oral, \en{I'm going to travel} devient souvent :}{\en{I'm wanna travel.}}{\en{I'm gotta travel.}}{\en{I'm gonna travel.}}
\qcm[b]{« des vacances », en américain, se dit :}{\en{holidays}}{\en{a vacation}}{\en{a congé}}
\qcm[a]{La Route 66 relie :}{Chicago et Los Angeles}{New York et Boston}{Miami et Seattle}
\qcm[c]{Un \en{mile} fait environ :}{500 mètres}{1 kilomètre}{1,6 kilomètre}
\end{testexpress}

\noindent{\fontsize{9.6}{12}\selectfont\color{otMuted}Niveau visé (CECRL, niveau A1) : « peut parler
de ses projets en termes simples ». Les corrigés de l'unité sont page~\pageref{cor-u16-1} ; les textes
des pistes d'exercices, page~\pageref{trans-16.10}.\par}
